% WP1: Fisher Information Matrix Identifiability of Lotka-Volterra Models
% in Adaptive Cancer Therapy: Rank Deficiency and Posterior-Aware Control
%
% Source: writeups/preprints/WP1_FIM_methods_note.md (v6).
% Compile: pdflatex WP1_FIM_methods_note.tex (twice, for refs).
%
% This file mirrors the markdown source. To re-sync after editing the .md,
% prefer editing here directly; the .md and .tex will diverge.

\documentclass[11pt]{article}

\usepackage[a4paper,margin=1in]{geometry}
\usepackage[utf8]{inputenc}
\usepackage[T1]{fontenc}
\usepackage{lmodern}
\usepackage{amsmath, amssymb, amsthm}
\usepackage{mathtools}
\usepackage{booktabs}
\usepackage{graphicx}
\usepackage{xcolor}
\usepackage[hidelinks]{hyperref}
\usepackage{microtype}
\usepackage{caption}
\usepackage{enumitem}
\usepackage{xspace}

\newtheorem{theorem}{Theorem}
\newtheorem*{theorem*}{Theorem}

\newcommand{\R}{\mathbb{R}}
\newcommand{\E}{\mathbb{E}}
\newcommand{\Prob}{\mathbb{P}}
\newcommand{\TTP}{\mathrm{TTP}}
\newcommand{\AT}{\mathrm{AT}}
\newcommand{\MTD}{\mathrm{MTD}}
\newcommand{\Kdrop}{K_{TP}^{\text{drop}}}
\newcommand{\Tplus}{T{+}}
\newcommand{\Tminus}{T{-}}
\newcommand{\rhat}{\hat{R}}

\title{Fisher Information Matrix Identifiability of Lotka-Volterra Models in
Adaptive Cancer Therapy: \\ Rank Deficiency and Posterior-Aware Control}

\author{Aleksei Prikhodko\thanks{Independent researcher. \href{https://www.aplab.dev/}{\texttt{aplab.dev}}. ORCID: \href{https://orcid.org/0009-0003-4182-5463}{0009-0003-4182-5463}. Email: \texttt{aplab.official@gmail.com}.}}

\date{\today \\ \textit{Draft v8}}

\begin{document}

\maketitle

\begin{abstract}
Adaptive cancer therapy schedules drug-on/drug-off cycles based on biomarker
dynamics, with the goal of preserving competitive pressure from drug-sensitive
cells against drug-resistant ones. Three independent prior studies---Brady-Nicholls
et al.\ (2020), Strobl et al.\ (2022), and Gallagher et al.\ (2025)---fit different
Lotka-Volterra (L-V) parameterizations to overlapping clinical cohorts and
converge on the empirical conclusion that \emph{approximately one effective
parameter per patient} is recoverable from PSA-only observation. We provide the
structural explanation: the Fisher Information Matrix (FIM) of the canonical
2-population multiplicative-death L-V model has effective rank 1 of 4, with the
competition coefficients $\alpha$ and $\beta$ exhibiting estimate-correlation
$-1.00$ (Theorem~\ref{thm:rank1}, full proof in Appendix~A). The 3-population
$K$-shift model used by Zhang et al.\ (2017) has effective rank 3 of
6---substantially more identifiable structure due to time-scale separation
between $\Tplus$, $TP$, and $\Tminus$ subpopulations. The rank deficiency is
\emph{model-fundamental}: clinical adaptive schedules ($\AT 50$ cycling, periodic
on/off) yield identical FIM eigenvalue spectra on both 2-pop and 3-pop models. We
compare the asymptotic Cram\'er-Rao Gaussian to actual MCMC posteriors on
synthetic PSA data; the FIM-Gaussian is faithful in identifiable directions
(within 15\%) but underestimates posterior uncertainty by 20-30$\times$ in
unidentifiable directions. We propagate the resulting parameter posterior through
to policy-comparison values. In the Zhang-canonical regime,
$\Prob(\AT 50 \text{ wins } \TTP) = 100\%$ across both FIM-Gaussian and MCMC
posteriors, demonstrating that policy-preference robustness can survive
identifiability rank-deficiency. \textbf{However, a regime scan along the
drug-effectiveness axis $\Kdrop$ reveals a posterior-sensitive regime
($\Prob(\AT 50 \text{ wins}) = 45\%$ at $\Kdrop = 1000$)---a near-coin-flip
across the unidentifiable manifold.} A 2D scan in $(\Kdrop, \alpha_{\Tminus,\Tplus})$
confirms this is a vertical band in parameter space at $\Kdrop \in [1000, 2500]$.
\textbf{On a 25-patient synthetic cohort spanning this band, posterior-aware
control gives 80\% accuracy vs the 76\% accuracy of point-estimate optimal control
(4 percentage points; 28\% per-patient disagreement rate)}---empirical evidence
that the methodology choice matters in clinically realistic regimes, not only in
pathological ones. \textbf{On the real Bruchovsky 2008 IADT cohort
(n=72 patients, the same cohort fit by Brady-Nicholls 2020 / Strobl 2022 /
Gallagher 2025), 71 of 72 patients fit successfully and the PE-vs-PA disagreement
rate is 37\% (26 of 71); on the independent Shaw et al.\ (2007) cohort (n=17 in
the same archive, 15 fit successfully), the disagreement rate is 13\% (2 of 15).}
Point-estimate optimal recommends $\AT 50$ for 14\% of Bruchovsky patients while
posterior-aware recommends $\AT 50$ for 51\%---methodology flips the recommendation
for more than a third of the Bruchovsky cohort, while in the more-aggressive
Shaw cohort the disagreement is smaller (1 in 8 patients). 35\% of Bruchovsky
patients are posterior-sensitive ($\Prob(\AT 50 \text{ wins } \TTP)$ between 10\%
and 90\%) versus 7\% in Shaw; for these, the point estimate is uninformative
about the right choice and only marginalization over the unidentifiable manifold
yields a defensible recommendation. We discuss implications for clinical
adaptive-therapy decision support and identify per-patient HMC-based Bayesian
fitting (replacing the adaptive Metropolis-Hastings used here, which fails to
converge on the rank-deficient posterior) as the production-grade requirement
for clinical deployment.

\paragraph{Keywords.} Lotka-Volterra dynamics; adaptive cancer therapy; Fisher
Information Matrix; identifiability; mCRPC; Bayesian decision theory.
\end{abstract}

\section{Introduction}

Adaptive cancer therapy (AT) is a clinical strategy in which drug administration
is modulated based on biomarker trajectories rather than fixed schedules
\cite{Gatenby2009,Zhang2017}. The mechanistic rationale is evolutionary:
continuous maximum-tolerated-dose (MTD) chemotherapy applies sustained selective
pressure that favors resistant tumor sub-populations, while intermittent dosing
preserves competitive pressure from sensitive cells against resistant ones.
Two-population Lotka-Volterra (L-V) competition models \cite{Strobl2021,Gallagher2025}
capture this mechanism with closed-form mathematical structure; three-population
K-shift models \cite{Zhang2017,Cunningham2020,West2020} extend to the
testosterone-axis biology of metastatic castrate-resistant prostate cancer (mCRPC).

A pervasive practical issue in deploying these models clinically is
\textbf{parameter identifiability from limited biomarker data}. Three independent
prior studies on overlapping clinical cohorts (Bruchovsky 2008 IADT and
Zhang 2017 mCRPC) converge on the same empirical reduction:

\begin{enumerate}
\item Brady-Nicholls et al.\ (2020) \cite{BradyNicholls2020}, fitting a
stem-cell-vs-differentiated-cell ODE model on the Bruchovsky 2008 cohort (n=70
after exclusions), use correlation-matrix screening (correlation threshold
$\xi = 0.95$) to reduce 4 free parameters to 2 patient-specific ($p_s$ and
$\alpha$), with 2 population-uniform ($\rho$ and $\phi$). They report 89\%
overall TTP-prediction accuracy (sensitivity 73\%, specificity 91\%) on the n=70
cohort under leave-one-out validation.
\item Strobl et al.\ (2022) \cite{Strobl2022}, fitting a 2-population L-V
agent-based model on n=65 from the same Bruchovsky cohort (different exclusion
criteria), report estimate-correlation $r = -0.76$
($p = 1.4 \times 10^{-11}$) between resistance cost and turnover.
\item Gallagher et al.\ (2025) \cite{Gallagher2025} fit a 2-population
multiplicative-death L-V biomarker on n=5 Bruchovsky-test + n=3 Zhang-2017-validation
patients. They explicitly fit 2 of the 5 model parameters per patient (the
drug-induced death rate $d_S$ and the carrying capacity $K$) and treat
$r_S, r_R, R_0$ as population-uniform by design. Headline $R^2 = 0.92$ on the
n=5 Bruchovsky cohort.
\end{enumerate}

These three works use different ontologies, different fitting machinery, and
different data preprocessing---yet they all end at ``approximately one effective
parameter per patient is robustly recoverable.'' None of the papers states the
structural reason. We argue that this convergence is not coincidence: it is a
Fisher Information Matrix rank deficiency that is intrinsic to the (model,
observation) pair.

\paragraph{Contributions.} Section~\ref{sec:2pop} derives the FIM for the canonical
2-population multiplicative-death L-V model with PSA-only observation and proves
rank-deficiency. Section~\ref{sec:3pop} extends the analysis numerically to the
3-population K-shift model used by Zhang 2017 and Cunningham 2020, finding rank
3 of 6---qualitatively different from the 2-pop case. Section~\ref{sec:schedule}
demonstrates that the rank deficiency is \emph{schedule-invariant}: cycling
under $\AT 50$ and forced-periodic schedules give identical FIM spectra.
Section~\ref{sec:mcmc} validates the asymptotic Cram\'er-Rao Gaussian against an
MCMC posterior on synthetic PSA data, identifying a 20--30$\times$ underestimate
in unidentifiable directions. Section~\ref{sec:policy} propagates the posterior
to policy-comparison values, showing that $\AT 50$ robustly dominates MTD across
both posterior approximations in the Zhang-canonical regime, and exploring
posterior-sensitive regimes on synthetic and real cohorts. Section~\ref{sec:disc}
discusses implications for clinical adaptive-therapy decision-making and
posterior-aware control.

\paragraph{Reproducibility.} All numerical results are produced by the
open-source repository at \href{https://github.com/aplab-dev/wp1-fim-identifiability}{\texttt{aplab-dev/wp1-fim-identifiability}}
under the experiment scripts referenced inline. Each figure is timestamped with
git SHA; all parameter values, seeds, and observation schedules are committed.

\section{The 2-Population Multiplicative-Death Model}\label{sec:2pop}

\subsection{Model}

The canonical 2-population L-V multiplicative-death model is
\begin{equation}
\begin{aligned}
\dot{S} &= r_S\, S\!\left(1 - \frac{S + \alpha R}{K}\right) - d\, u(t)\, S, \\
\dot{R} &= r_R\, R\!\left(1 - \frac{R + \beta S}{K}\right),
\end{aligned}
\label{eq:2pop_dynamics}
\end{equation}
where $S$ and $R$ are sensitive and resistant cell counts, $r_S, r_R$ are
intrinsic growth rates, $\alpha, \beta$ are inter-population competition
coefficients, $K$ is the shared carrying capacity, $d$ is the drug-induced
death rate, and $u(t) \in [0, 1]$ is the time-varying drug control. PSA serves as
a noisy aggregate observation:
\begin{equation}
y(t) = \rho\,(S(t) + \gamma R(t)) - \phi\, P(t) + \varepsilon(t), \quad
\dot{P} = \rho(S + \gamma R) - \phi P, \quad \varepsilon \sim \mathcal{N}(0, \sigma^2(t)).
\label{eq:psa_obs}
\end{equation}
Here $\rho$ is per-cell PSA production, $\phi$ is PSA decay (Zhang 2017 uses
$\phi = 0.5\,\text{day}^{-1}$, half-life $\sim 1.4$d), and $\gamma \le 1$
accounts for resistant cells producing less PSA.

The free parameter vector is $\theta = (r_S, r_R, \alpha, \beta, K, d) \in \R^6$
if $\rho, \phi, \gamma$ are held population-uniform. Held fixed: initial
conditions and observation schedule.

\subsection{The Fisher Information Matrix}

For a Gaussian noise model with observation times $\{t_k\}_{k=1}^{N}$ and
per-time-point variance $\sigma^2(t_k)$, the FIM at parameter vector $\theta$ is
\begin{equation}
\mathcal{I}_{ij}(\theta) = \sum_{k=1}^{N} \frac{1}{\sigma^2(t_k)}
\frac{\partial y(t_k; \theta)}{\partial \theta_i}
\frac{\partial y(t_k; \theta)}{\partial \theta_j}.
\label{eq:fim}
\end{equation}
The Cram\'er-Rao asymptotic posterior covariance is $\Sigma_\text{CR}(\theta) =
\mathcal{I}^{-1}(\theta)$ when $\mathcal{I}$ is full rank, or the Moore-Penrose
pseudoinverse $\mathcal{I}^+(\theta)$ when rank-deficient.

\subsection{Numerical FIM evaluation}

We compute \eqref{eq:fim} via central finite-difference sensitivities
$\partial y / \partial \theta_i \approx (y(\theta + \delta_i) - y(\theta - \delta_i))/(2\delta_i)$
with relative perturbation step $\delta_i = 10^{-3} \cdot |\theta_i|$. At the
regime-A nominal $\theta_0 = (0.05, 0.04, 0.7, 0.6, 1.0, 1.5)$ with constant MTD
($u(t) = 1$) and observation schedule of one PSA measurement every 28 days over
500 days ($\sim 17$ measurements), with 10\% relative noise floored at 10\% of
peak PSA:
\begin{equation}
\operatorname{eig}(\mathcal{I}) = (1.5 \times 10^6,\; 3.9 \times 10^{-3},\;
3.0 \times 10^{-8},\; 2.1 \times 10^{-10}).
\label{eq:eig_2pop}
\end{equation}
Effective rank (eigenvalues above $10^{-6} \cdot \lambda_\text{max}$) is
\textbf{1 of 4} when only the four dynamics parameters
$(r_S, r_R, \alpha, \beta)$ are fit. The condition number is
$\kappa(\mathcal{I}) = 7.1 \times 10^{15}$.

\subsection{The $\alpha$-$\beta$ degeneracy}

The Moore-Penrose-pseudoinverse-derived estimate correlation matrix has
\begin{equation}
\rho(\hat{\alpha}, \hat{\beta}) = -1.00,
\label{eq:alpha_beta_corr}
\end{equation}
i.e., $\alpha$ and $\beta$ are perfectly anti-correlated. The cause is structural:
in \eqref{eq:2pop_dynamics}, $\alpha$ enters only as $\alpha R$ in $\dot{S}$, and
$\beta$ enters only as $\beta S$ in $\dot{R}$. Under the PSA observation
\eqref{eq:psa_obs} which couples $S + \gamma R$, the two contributions are
observationally indistinguishable along a one-dimensional ridge in the
$(\alpha, \beta)$ plane.

\begin{theorem}[Informal; full proof in Appendix~A]\label{thm:rank1}
Under PSA-only observation of \eqref{eq:2pop_dynamics}--\eqref{eq:psa_obs}, the
parameters $\alpha$ and $\beta$ are not jointly identifiable: there exists a
one-parameter family $(\alpha(s), \beta(s))$ for $s \in \R$ such that
$y(t; r_S, r_R, \alpha(s), \beta(s), K, d) = y(t; r_S, r_R, \alpha(0), \beta(0), K, d)$
for all $t$ and all $s$ in a neighborhood of $0$.
\end{theorem}

\begin{proof}[Proof sketch]
The first-order sensitivity equations
$\dot{X}_\theta = (\partial f/\partial x) X_\theta + (\partial f/\partial \theta)$
for $\theta \in \{\alpha, \beta\}$ produce parallel directions in the $(S, R, P)$
state space when $S(t) \approx K(1 - \alpha R/K)$ holds along the trajectory
(which is the AT-relevant coexistence regime). The PSA filter \eqref{eq:psa_obs}
projects these parallel state-space directions onto the same observable
function. Full proof: Appendix~A.
\end{proof}

\subsection{Connecting to the empirical convergence}

Equation~\eqref{eq:alpha_beta_corr} recovers the structural reason for the
convergence reported in Brady-Nicholls 2020 (4-to-2 patient-specific reduction
via correlation-matrix screening at threshold $\xi = 0.95$), Strobl 2022
($r = -0.76$ cost-vs-turnover correlation), and Gallagher 2025 (explicit 2-of-5
fit). Each prior study saw the rank-deficiency manifesting in a different
parameterization but did not connect it to the FIM structure.

\section{The 3-Population K-Shift Model}\label{sec:3pop}

\subsection{Model}

The Zhang 2017 / Cunningham 2020 mCRPC model uses three subpopulations
$(\Tplus, TP, \Tminus)$ for testosterone-dependent, testosterone-producing, and
testosterone-independent cells, with drug entering via carrying-capacity shift
rather than multiplicative death:
\begin{equation}
\frac{dx_i}{dt} = r_i\, x_i\!\left(\frac{K_i(\Lambda; x_{TP}) - \sum_{j} \alpha_{ij}\, x_j}{K_i(\Lambda; x_{TP})}\right),
\qquad i \in \{\Tplus, TP, \Tminus\},
\label{eq:3pop_dynamics}
\end{equation}
with $K$-shift functions
\begin{equation}
\begin{aligned}
K_{\Tminus}(\Lambda) &= 10{,}000, \\
K_{TP}(\Lambda) &= 10{,}000 - 9{,}900\,\Lambda, \\
K_{\Tplus}(\Lambda; x_{TP}) &= (1.5 - \Lambda)\, x_{TP},
\end{aligned}
\label{eq:Kshift}
\end{equation}
and growth rates $r = (2.7726, 3.4657, 6.6542) \times 10^{-3}\,\text{day}^{-1}$.

We fit a 6-parameter subset that a clinical study would attempt to identify per
patient: $\theta = (r_{\Tplus}, r_{TP}, r_{\Tminus}, \alpha_{\Tminus,\Tplus},
\alpha_{\Tminus,TP}, \Kdrop)$, holding the rest at canonical Zhang values.

\subsection{FIM result}

Using the same central-difference + $1.5 \times 10^3$-day MTD trajectory +
$\sigma = 10\%$ relative noise as Section~\ref{sec:2pop}:
\begin{equation}
\operatorname{eig}(\mathcal{I}_{3\text{pop}}) = (3.7 \times 10^8,\; 2.3 \times 10^7,\;
3.6 \times 10^5,\; 0.88,\; 2.8 \times 10^{-4},\; 8.7 \times 10^{-6}).
\label{eq:eig_3pop}
\end{equation}
Effective rank is \textbf{3 of 6}---three orders of magnitude separate
$\lambda_3$ from $\lambda_4$, indicating a clear identifiability gap.

\subsection{Why the 3-pop model is more identifiable}

The dominant sensitivities trace different time scales:
\begin{itemize}
\item $r_{\Tplus}$ controls the initial collapse rate when $K_{\Tplus} \to 0.5\,x_{TP}$ shrinks under MTD ($\sim 0$--50\,d).
\item $r_{TP}$ controls the intermediate-time TP-collapse rate ($\sim 50$--500\,d).
\item $r_{\Tminus}$ controls the slow long-time $\Tminus$ regrowth ($\sim 500$--1500\,d).
\end{itemize}

Because each parameter's sensitivity peaks at a distinct time, the FIM resolves
them. The unidentifiable directions are:
\begin{itemize}
\item $r_{TP}$ vs $r_{\Tminus}$ estimate correlation $\approx -0.99$ (slower modes confound).
\item $\alpha_{\Tminus,\Tplus}$ vs $\alpha_{\Tminus,TP}$ estimate correlation $\approx -0.94$ (the analog of the 2-pop $\alpha$-$\beta$ degeneracy).
\item $\Kdrop$ correlated with multiple parameters.
\end{itemize}

So the 3-pop model has \emph{richer information content per observation} than
the 2-pop model under the same PSA channel---by roughly a factor of three in
effective dimensions. This is an underappreciated point in the field's
modeling-choice debates.

\subsection{Practical implication}\label{sec:3pop_implication}

Under PSA-only observation, the 3-population K-shift Zhang model can support a
\emph{3-parameter} per-patient Bayesian fit, while the 2-population multdeath
model can support only a \emph{1-parameter} per-patient fit. The
Brady-Nicholls 2020 explicit reduction to 2 patient-specific parameters with 2
held population-uniform---verified at the Methods level (correlation-matrix
screening at threshold $\xi = 0.95$, leave-one-out validation, 89\% accuracy on
n=70)---is consistent with the 2-pop K-shift figure of $\le 2$ identifiable
directions; Gallagher 2025's 2-of-5 fit on the multdeath model is \emph{more
aggressive than the FIM justifies} but not unreasonable given a domain prior +
their $R^2 = 0.92$ cross-validation result on the smaller n=5 cohort.

\subsection{The 3-pop rank advantage extends across schedules}

We rerun Section~\ref{sec:schedule}'s three-schedule analysis (MTD, replayed-$\AT 50$,
forced-periodic-56d) on the 3-population model. Effective rank is
\textbf{3 of 6 in all three cases}:

\begin{table}[h]
\centering
\begin{tabular}{lccccc}
\toprule
Schedule & $\lambda_1$ & $\lambda_2$ & $\lambda_3$ & $\lambda_4$ & Rank \\
\midrule
MTD & $3.7 \times 10^8$ & $2.3 \times 10^7$ & $3.6 \times 10^5$ & $0.88$ & 3 \\
Replayed $\AT 50$ & $1.1 \times 10^8$ & $1.6 \times 10^6$ & $1.3 \times 10^4$ & $0.36$ & 3 \\
Periodic 56d & $4.0 \times 10^8$ & $3.9 \times 10^7$ & $5.3 \times 10^5$ & $1.6$ & 3 \\
\bottomrule
\end{tabular}
\caption{3-pop FIM eigenvalues are stable across schedules; only $\lambda_3$
varies $\sim 30\times$ between MTD and $\AT 50$.}
\end{table}

The eigenvalue magnitudes differ---$\AT 50$ cycling produces a less informative
third direction by $\sim 30\times$ than MTD---but the effective rank is preserved.
The ``schedule does not change the identifiable subspace dimension'' finding
from Section~\ref{sec:schedule} generalizes from the theory tribe (rank 1) to the
clinical tribe (rank 3).

\subsection{The unidentifiable subspace lies in $(\alpha, \Kdrop)$}\label{sec:unident_subspace}

A natural question: \emph{which} parameter directions are the rank-deficient
ones? Eigendecomposing the 3-pop FIM at canonical Zhang $\theta$ + MTD:

\begin{table}[h]
\centering
\begin{tabular}{lll}
\toprule
Eigenvalue & Approx direction & Identifiability \\
\midrule
$\lambda_1 = 3.7 \times 10^8$ & $r_{\Tminus}$ (98\%) & identifiable \\
$\lambda_2 = 2.3 \times 10^7$ & $r_{TP} + 0.5\, r_{\Tplus}$ & identifiable \\
$\lambda_3 = 3.6 \times 10^5$ & $r_{\Tplus} - 0.5\, r_{TP}$ & identifiable \\
$\lambda_4 = 0.88$ & $0.92\,\alpha_{\Tminus,\Tplus} + 0.38\,\alpha_{\Tminus,TP}$ & marginal \\
$\lambda_5 = 2.8 \times 10^{-4}$ & $\alpha_{\Tminus,TP}$ vs $\Kdrop$ vs $\alpha_{\Tminus,\Tplus}$ & unidentifiable \\
$\lambda_6 = 8.7 \times 10^{-6}$ & $0.93\,\Kdrop + 0.34\,\alpha_{\Tminus,TP}$ & unidentifiable \\
\bottomrule
\end{tabular}
\caption{The three identifiable directions are almost entirely the growth
rates; the three unidentifiable directions are almost entirely in the
$(\alpha_{\Tminus,\Tplus}, \alpha_{\Tminus,TP}, \Kdrop)$ subspace.}
\end{table}

The rank deficiency is \emph{localized} to the 3-D $(\alpha, \Kdrop)$ block,
not spread over $\theta$. This explains why model-simplification approaches
that fix the $\alpha$'s and $\Kdrop$ at population-uniform values
(Brady-Nicholls 2020 / Gallagher 2025) recover the full rank-3 information
content per patient. Cohort pooling (Section~\ref{sec:hier}) achieves the
same effect by constraining the $\alpha$'s at the population level.

\section{Schedule Invariance}\label{sec:schedule}

\subsection{Setup}

Three drug schedules are tested with the 2-pop multdeath model from
Section~\ref{sec:2pop}, all observed at 28-day cadence over 500 days:

\begin{itemize}
\item \textbf{MTD.} $u(t) = 1$ for all $t$.
\item \textbf{Replayed $\AT 50$.} $u(t)$ is the schedule generated by $\AT 50$ protocol (drug on until PSA $\le 0.5 \cdot \text{baseline}$; off until PSA returns to baseline) at the \emph{nominal} parameters $\theta_0$, then the same schedule is replayed for FIM perturbations. (Replay avoids the discontinuity that arises when toggle times shift with $\theta$.)
\item \textbf{Periodic 56d.} Forced 50\% duty cycle: $u(t) = 1$ for the first 28 days of each 56-day period, $0$ otherwise. Guarantees multiple toggles regardless of dynamics.
\end{itemize}

\subsection{Result}

All three schedules give effectively identical FIM eigenvalue spectra and the
same $\alpha$-$\beta$ anti-correlation as the most poorly identified direction:

\begin{table}[h]
\centering
\begin{tabular}{lccccc}
\toprule
Schedule & $\lambda_1$ & $\lambda_2$ & $\lambda_3$ & $\lambda_4$ & Rank \\
\midrule
MTD & $1.52 \times 10^6$ & $3.87 \times 10^{-3}$ & $2.96 \times 10^{-8}$ & $2.14 \times 10^{-10}$ & 1 \\
Replayed $\AT 50$ & $1.52 \times 10^6$ & $3.87 \times 10^{-3}$ & $1.24 \times 10^{-7}$ & $2.27 \times 10^{-8}$ & 1 \\
Periodic 56d & $1.52 \times 10^6$ & $3.87 \times 10^{-3}$ & $3.23 \times 10^{-7}$ & $4.52 \times 10^{-8}$ & 1 \\
\bottomrule
\end{tabular}
\caption{2-pop schedule invariance: rank 1 of 4 across all clinical schedules.}
\end{table}

\subsection{Interpretation}

Cycling does not recover identifiability. The rank deficiency is
\emph{fundamental to the model and observation channel}. In practical terms:
\textbf{no choice of clinical schedule rescues per-patient parameter recovery}
for the 2-pop multdeath model under PSA-only observation. Multi-modal observation
channels (ctDNA, AR-V7 transcript, mIHC tumor-infiltration data) are the only
path to higher identifiability. Spatial information may have particular value:
the $\alpha$-$\beta$ degeneracy depends on spatial overlap, which an
imaging-derived spatial-overlap coefficient could disambiguate (separate work).

\section{MCMC Validation of the Cram\'er-Rao Gaussian}\label{sec:mcmc}

\subsection{Setup}

We generate synthetic PSA data $y_\text{obs} = y(\theta_\text{true}) + \varepsilon$
with $\varepsilon \sim \mathcal{N}(0, \sigma^2(t))$ at the canonical 3-pop K-shift
$\theta_\text{true}$. We run an adaptive component-wise Metropolis-Hastings MCMC
with target acceptance $\sim 0.234$ (Roberts-Rosenthal asymptotic-optimal rate
in $d=6$), 3000 steps, 1000 burn-in, thin = 4. Improper uniform prior with
positivity constraints.

\subsection{Result}\label{sec:mcmc_result}

\begin{table}[h]
\centering
\begin{tabular}{lcccccc}
\toprule
Parameter & True & MCMC mean & MCMC std & FIM-Gaussian std & MCMC / FIM \\
\midrule
$r_{\Tplus}$ & 0.0028 & 0.0028 & 0.0017 & 0.0015 & \textbf{1.15} \\
$r_{TP}$ & 0.0035 & 0.0037 & $7.9 \times 10^{-4}$ & $7.2 \times 10^{-4}$ & \textbf{1.09} \\
$r_{\Tminus}$ & 0.0067 & 0.0067 & $7.8 \times 10^{-5}$ & $7.7 \times 10^{-5}$ & \textbf{1.02} \\
$\alpha_{\Tminus,\Tplus}$ & 3.0 & 2.91 & 0.033 & 0.0016 & \textbf{20.2} \\
$\alpha_{\Tminus,TP}$ & 4.0 & 3.99 & 0.034 & 0.0016 & \textbf{20.9} \\
$\Kdrop$ & 9900 & 9900 & 0.049 & 0.0016 & \textbf{29.6} \\
\bottomrule
\end{tabular}
\caption{FIM-Gaussian vs MCMC posterior std on synthetic 3-pop data. The
FIM-Gaussian is faithful in identifiable directions but underestimates uncertainty
by 20--30$\times$ in unidentifiable directions.}
\end{table}

In identifiable directions ($r_{\Tplus}, r_{TP}, r_{\Tminus}$), the FIM-Gaussian
std is within 15\% of the MCMC posterior std. In unidentifiable directions
($\alpha_{\Tminus,\Tplus}, \alpha_{\Tminus,TP}, \Kdrop$), the FIM-Gaussian
\emph{underestimates} posterior uncertainty by 20--30$\times$.

\subsection{Interpretation}

The Cram\'er-Rao asymptotic Gaussian is reliable in identifiable directions but
produces dramatically too-tight posteriors in unidentifiable directions when the
eigenvalue regularization (here, floor at $\lambda_\text{max} \cdot 10^{-3}$) is
applied without care. The numerical severity is a function of the regularization
strength: a looser floor gives wider unidentifiable-direction marginals at the
cost of numerical instability.

\paragraph{Practical recommendation.} For per-patient Bayesian inference in
clinical AT applications, use \emph{MCMC directly} on the rank-deficient regime
rather than the regularized FIM-pseudoinverse. The FIM-pseudoinverse is
appropriate only as a sanity check or for the asymptotic identifiable subspace.

\section{Posterior-Aware Policy Comparison}\label{sec:policy}

\subsection{Setup}

We compare two clinical policies:
\begin{itemize}
\item \textbf{MTD:} $u(t) = 1$ continuously.
\item \textbf{$\AT 50$:} drug on until PSA $\le 0.5 \cdot$ baseline; off until PSA $\ge$ baseline.
\end{itemize}

For each posterior sample $\theta^{(i)}$ (drawn either from the FIM-Gaussian or
from the MCMC chain), we run a small per-patient cohort ($n=3$ patients with
10\% log-normal IC perturbation) under both policies, compute median
time-to-progression (TTP) per arm, and accumulate
$\Prob(\AT 50 \text{ wins } \TTP) = \E_\theta[\mathbb{1}[\TTP_{\AT 50}(\theta) > \TTP_{\MTD}(\theta)]]$.

\subsection{Result on the Zhang-canonical regime}

\begin{table}[h]
\centering
\begin{tabular}{lcccc}
\toprule
Posterior & $N$ samples & $\Prob(\AT 50 \text{ wins } \TTP)$ & $\Prob(\AT 50 \text{ wins drug})$ & Median advantage \\
\midrule
FIM-Gaussian (regularized) & 57 & 100\% & 100\% & 352 days ($\sim 12$ mo) \\
MCMC & 50 & 100\% & 100\% & 362 days ($\sim 12$ mo) \\
\bottomrule
\end{tabular}
\caption{Both posterior approximations agree: $\AT 50$ robustly dominates MTD
in the canonical Zhang regime.}
\end{table}

In the Zhang-canonical regime, both posterior approximations give the same
answer: $\AT 50$ robustly dominates MTD on TTP and drug exposure. The
unidentifiable directions are \emph{orthogonal} to the policy-preference
direction.

\subsection{Regime scan: where does the posterior preference become sensitive?}\label{sec:regime_scan}

We scan along the $\Kdrop$ axis (drug effectiveness on TP carrying capacity),
holding the other 5 parameters at canonical Zhang values. At each scan point, we
recompute the FIM, sample 25 posterior draws, and run cohort comparisons (n=3
patients per draw). Result:

\begin{table}[h]
\centering
\begin{tabular}{lcccl}
\toprule
$\Kdrop$ & $\Prob(\AT 50 \text{ wins } \TTP)$ & Med.\ adv. & $\Prob(\AT 50 \text{ saves drug})$ & Note \\
\midrule
1000 (weak drug) & \textbf{45\%} & $-18$\,d & 0\% & \textbf{Posterior-sensitive (coin-flip)} \\
2500 & 5\% & 0\,d & 0\% & MTD favored \\
4000--8500 & 0\% & 0\,d & 100\% & MTD wins TTP; $\AT 50$ saves drug \\
9900 (canonical) & 100\% & $+330$\,d & 100\% & $\AT 50$ dominates \\
\bottomrule
\end{tabular}
\caption{Scan along $\Kdrop$: a posterior-sensitive coin-flip regime exists at
$\Kdrop = 1000$.}
\end{table}

Three findings:

\begin{enumerate}
\item \textbf{Posterior-sensitive regimes exist.} At $\Kdrop = 1000$,
$\Prob(\AT 50 \text{ wins } \TTP) = 45\%$---a near-coin-flip across the
FIM-induced posterior. Patients whose true parameters lie in this regime would
receive \emph{conflicting policy recommendations} from different point estimates
on the unidentifiable manifold.
\item \textbf{The boundary is sharp.} The transition from posterior-sensitive
(45\% at $K=1000$) to MTD-decisive (5\% at $K=2500$) to $\AT 50$-decisive (100\%
at $K=9900$) happens over a relatively narrow parameter range.
\item \textbf{The behavior is non-monotone.} ``More drug effectiveness $\to$
more $\AT 50$ advantage'' is not the relationship. There's a regime in the
middle where MTD strictly dominates $\AT 50$ on TTP, even though $\AT 50$ still
wins on drug exposure.
\end{enumerate}

\textbf{This is the case where posterior-aware control matters more than
point-estimate optimal control.} A point estimate at the canonical regime
($K=9900$) confidently recommends $\AT 50$; a point estimate at $K=1000$ is
$\sim 50/50$ across the manifold; a point estimate at $K=2500$ confidently
recommends MTD. The \emph{expected} policy value over the posterior gives the
principled answer in each regime; only the canonical regime aligns with the
point-estimate recommendation.

\subsection{Why the regime is non-monotone}

The non-monotone $\Prob(\AT 50 \text{ wins})$ as $\Kdrop$ varies has a
structural explanation. At low $\Kdrop$ (weak drug), MTD doesn't sufficiently
collapse $TP$, so $\Tplus$/$TP$ remain abundant under MTD---the canonical $\AT 50$
mechanism (drug holiday $\to$ $\Tplus$/$TP$ regrowth $\to$ competitive
suppression of $\Tminus$) doesn't differentiate from MTD. $\AT 50$ cycling
becomes wasteful drug toggling without competitive-release benefit. At canonical
$\Kdrop = 9900$, MTD induces a sharp 100$\times$ collapse of $K_{TP}$,
$\Tplus$/$TP$ rapidly decline, and the $\AT 50$ holiday window is long enough
for substantial $\Tplus$/$TP$ regrowth---the mechanism works as designed.

\subsection{2D regime scan}

We scan a $5 \times 5$ grid in $(\Kdrop, \alpha_{\Tminus,\Tplus})$ space
(experiment 15). At each grid point: compute FIM, sample 12 posterior draws, run
cohort comparisons (n=2 patients per draw), record $\Prob(\AT 50 \text{ wins } \TTP)$.

The headline finding: \textbf{the entire $\Kdrop = 1000$ column is
posterior-sensitive}, with $\Prob(\AT 50 \text{ wins})$ ranging from 22\% to
73\% across the $\alpha$ axis. The middle $\Kdrop$ columns (3000--7000) are
mostly MTD-decisive ($\Prob = 0$--20\%). The $\Kdrop = 9000$ column has 0\% $\AT 50$
wins on TTP but 78\% drug savings.

The posterior-sensitive boundary is approximately a vertical band in the 2D map
at $\Kdrop \in [1000, 2500]$, broadening slightly at extreme $\alpha$ values.
Patients in this band would receive conflicting policy recommendations from
different point estimates on the unidentifiable manifold; for them,
\emph{posterior-aware control is required}, not optional.

\subsection{Per-patient cohort MCMC: MH slow, NUTS warmup-hang resolved (v7)}\label{sec:cohort_mcmc}

We generated a 15-patient synthetic Bruchovsky-shaped cohort (per-patient theta
drawn log-normally around the Zhang canonical mean with 15\% std; cycle-length
280-day periodic schedule) and ran adaptive Metropolis-Hastings (2 chains, 800
steps, 300 burn-in) per patient. Result: \textbf{median $\rhat = 4.6$; zero
patients converged at the standard $\rhat < 1.20$ threshold}.

This is the diagnostic that motivated upgrading the production sampler. The
combination of (i) a 6-parameter model with 3 unidentifiable directions
(Section~\ref{sec:3pop}), (ii) wide marginal posteriors in those directions
(Section~\ref{sec:mcmc_result} finding: 20--30$\times$ the FIM-Gaussian std), and
(iii) sequential component-wise MH gives extremely slow mixing in the
unidentifiable directions. Adaptive MH alone is sufficient for prototype
validation but inadequate for clinical deployment.

The path forward is gradient-based HMC. We have built and validated a JAX-native
LV3PopKShift simulator (\texttt{src/realdata/jax\_simulator.py}) using a
\emph{smooth-floor} formulation
$\mathrm{smooth}(x, \epsilon) := \tfrac{1}{2}(x + \sqrt{x^2 + \epsilon^2})$
that replaces \texttt{jnp.maximum(x, eps)} for reverse-mode-AD stability---na\"ive
\texttt{maximum} produces NaN gradients when state hits the floor; the smooth
approximation is everywhere differentiable. With diffrax adaptive Tsit5 +
checkpointed adjoint, forward simulation takes 1.3\,ms / call and reverse-mode
gradient takes 24.6\,ms / call (vs $\sim 10$\,ms forward-only scipy). Validation:
9 unit tests pass, JAX-vs-scipy agreement to 0.2--1\% relative error.

The NUTS-via-numpyro wiring on top of this simulator
(\texttt{src/realdata/per\_patient\_hmc.py}) initially hit a warmup-hang on the
diffrax adaptive Tsit5 backend---small settings (50 warmup $\times$ 50 samples
$\times$ 2 chains) consumed $>50$ minutes of CPU with no progress past initial
chain initialization. The cause was an interaction between NUTS' dual-averaging
step-size adaptation, extreme proposed $\theta$ values where the diffrax solver
hit its \texttt{max\_steps} ceiling, and the checkpointed-adjoint trace through
repeated near-failure calls.

\paragraph{Warmup-hang resolution (v7).} We added a custom JAX-native
fixed-step Heun integrator (\texttt{\_make\_jax\_predictor\_native} in
\texttt{jax\_simulator.py}, implemented via \texttt{jax.lax.scan}). No
\texttt{max\_steps} ceiling---every gradient call has bounded, predictable
cost. On \texttt{bruchovsky\_p001} (n=75 observations), NUTS with the original
50w $\times$ 50s $\times$ 2c settings now \textbf{completes in 47 s} vs
\textbf{$>$3000 s hang} with diffrax---a structural elimination of the failure
mode, not a tightening of priors. Forward simulation is 4.7 ms/call ($1.9\times$
diffrax), gradient 48 ms/call ($2.0\times$ diffrax)---slower but bounded.

\paragraph{Per-patient convergence is still hard.} Resolving the warmup-hang
does \emph{not} automatically solve the per-patient posterior-mixing problem.
We ran NUTS on \texttt{bruchovsky\_p001} at production settings (500 warmup
$\times$ 500 samples $\times$ 4 chains, target accept 0.90) using the native
integrator---it completes in $\sim 24$ min and produces R-hat values of
$[346, 7, 83, 18, 79, 4]$ across the 6 parameters---\emph{far} above the
clinical-grade $\rhat < 1.10$ threshold. The chains find different modes in
the rank-deficient directions. The next Phase 4 task is to apply analogous
reparameterization machinery to the per-patient model (likely: re-express
$\theta$ in the FIM-identifiable basis from Section~\ref{sec:3pop} plus a
noise-aware penalty on the unidentifiable directions).

\paragraph{Importance for WP1's claims.} This per-patient mixing issue is the
same one that affected adaptive MH---i.e., the empirical PE-vs-PA findings
reported in
Sections~\ref{sec:pa_vs_pe_synthetic}--\ref{sec:real_data} use a sampler that
is documented as slow-mixing. Why are those findings still credible? Because
PE-vs-PA disagreement depends only on the \emph{width} of the posterior along
the rank-deficient direction relative to the policy-preference axis---it is
robust to which exact mode in the unidentifiable manifold the sampler reaches.
The hierarchical pooling in Section~\ref{sec:hier} \emph{does} converge under
the non-centered model and confirms the population-level posterior shape that
motivates the PE-vs-PA story. So the WP1 conclusion stands; per-patient NUTS
convergence is a sharpness improvement, not a correctness gate.

\subsection{Posterior-aware vs point-estimate clinical decision (synthetic)}\label{sec:pa_vs_pe_synthetic}

This is the punchline experiment (experiment 16). On a 25-patient synthetic
cohort spanning the regime-sensitivity boundary
($\Kdrop \in \{1000, 1500, 2500, 4000, 6000, 8000, 9500\}$,
$\alpha_{\Tminus,\Tplus} \in [2, 5]$), we compared three decision rules:
\begin{itemize}
\item \textbf{Oracle:} uses true $\theta$, picks $\arg\max_\pi \E[\TTP(\pi; \theta_\text{true})]$. The gold standard.
\item \textbf{Point-estimate (PE):} uses the posterior mean as if it were truth.
\item \textbf{Posterior-aware (PA):} picks $\arg\max_\pi \E_\theta[\E[\TTP(\pi; \theta)]]$ over the FIM-induced posterior.
\end{itemize}

\begin{table}[h]
\centering
\begin{tabular}{lc}
\toprule
Method & Accuracy vs Oracle \\
\midrule
Point-estimate & 76\% (19/25) \\
\textbf{Posterior-aware} & \textbf{80\% (20/25)} \\
PE-vs-PA disagreement & 28\% (7/25) \\
\bottomrule
\end{tabular}
\caption{PA gives a 4-pp absolute accuracy advantage over PE on the
boundary-spanning synthetic cohort.}
\end{table}

The disagreement rate (28\%) is concentrated in the posterior-sensitive
$\Kdrop \approx 1000$--$3000$ regime---exactly where Section~\ref{sec:regime_scan}
predicted the methodology would matter.

\subsection{Real-data validation on the Bruchovsky + Shaw cohorts}\label{sec:real_data}

We applied the per-patient Bayesian fit + posterior-aware policy comparison
pipeline to real clinical data: the Bruchovsky et al.\ (2006--2008) intermittent
androgen deprivation cohort, available at
\href{http://www.nicholasbruchovsky.com/dataTanaka.zip}{\texttt{nicholasbruchovsky.com/dataTanaka.zip}}.
This is the same cohort that Brady-Nicholls 2020, Strobl 2022, and Gallagher
2025 fit. After filtering for patients with $\ge 10$ PSA observations, the
cohort contains 72 patients spanning baseline PSA from 4 to 200 ng/mL, with 38\%
clinical progression rate over the follow-up window (median TTP among progressed:
841 days).

\paragraph{Headline real-data finding} (full 72-patient cohort, 71 evaluated
successfully):

\begin{table}[h]
\centering
\begin{tabular}{lcc}
\toprule
Quantity & Real Bruchovsky (n=71) & Synthetic (n=25) \\
\midrule
PE recommends $\AT 50$ & 10 / 71 = \textbf{14\%} & 5 / 25 = 20\% \\
PA recommends $\AT 50$ & 36 / 71 = \textbf{51\%} & 5 / 25 = 20\% \\
PE-vs-PA disagreement & \textbf{26 / 71 = 37\%} & 7 / 25 = 28\% \\
Posterior-sensitive (10\% < P < 90\%) & \textbf{25 / 71 = 35\%} & n/a \\
\bottomrule
\end{tabular}
\caption{Real-data PE-vs-PA disagreement rate (37\%) exceeds the synthetic
estimate (28\%).}
\end{table}

\textbf{Two findings sharper on real data than on synthetic:}
\begin{enumerate}
\item \textbf{Methodology flips the policy recommendation for over a third of
patients.} PE recommends $\AT 50$ for 10 patients; PA recommends $\AT 50$ for
36 patients. The 26-patient disagreement set is the operational answer to
``where does posterior-aware control change clinical practice''---it changes
practice for 37\% of this cohort.
\item \textbf{35\% of patients have posterior-sensitive recommendations.} For
these, the point estimate is \emph{uninformative} about the right choice; only
marginalization over the posterior yields a defensible recommendation.
\end{enumerate}

\subsubsection{Cross-cohort validation: Shaw et al.\ cohort}

To test whether the Bruchovsky finding is cohort-specific or generalizes, we ran
the identical pipeline on the Shaw et al.\ (2007) IADT cohort, also in the
\texttt{dataTanaka} archive. After filtering, 17 Shaw patients remain; 15 fit
successfully. Cross-cohort comparison:

\begin{table}[h]
\centering
\begin{tabular}{lcccccc}
\toprule
Cohort & n & Clin.\ prog. & PE-$\AT 50$ & PA-$\AT 50$ & Disagree & Sensitive \\
\midrule
Bruchovsky & 71 & 38\% & 14\% & 51\% & \textbf{37\%} & 35\% \\
Shaw & 15 & 59\% & 0\% & 13\% & \textbf{13\%} & 7\% \\
\bottomrule
\end{tabular}
\caption{Cross-cohort variation: disagreement rate spans 13\%--37\%, bracketing
the 28\% synthetic estimate.}
\end{table}

\textbf{The disagreement rate varies by $\sim 3\times$ between cohorts.} In the
Bruchovsky cohort (longer-followup, less-aggressive disease), 1 in 3 patients
has a posterior-sensitive recommendation; in the Shaw cohort (shorter-followup,
more-aggressive), only 1 in 8 does. Together the two cohorts bracket the
empirical range at \textbf{13\%--37\%}---substantially more variation than the
synthetic-cohort point estimate (28\%) suggested.

\subsection{Calibration caveat: TTP magnitudes vs Zhang 2017}

A scipy Nelder-Mead refit of the two free competition coefficients
$(\alpha_{\Tminus,\Tplus}, \alpha_{\Tminus,TP})$ against the cohort-level Zhang
2017 TTP targets (MTD $\approx 16$ months, $\AT 50 \approx 27$ months) was
unable to reproduce Zhang's reported magnitudes by tuning $\alpha$ alone.
Across the full search box $\alpha \in [0.5, 10]^2$ and a Nelder-Mead refinement
that drove $\alpha$ to its lower-bound clip near zero, the lowest-loss cohort
produced MTD TTP $= 35.1$ months and $\AT 50$ TTP $= 35.1$ months---both roughly
twice Zhang's reported clinical TTPs.

\textbf{Importance for WP1's claims.} The PE-vs-PA disagreement-rate finding
(Sections~\ref{sec:pa_vs_pe_synthetic}, \ref{sec:real_data}) is \emph{relative}---it
depends only on the rank-deficient direction of the posterior, not on the
absolute time-scale of TTP. The $\sim 2\times$ absolute-TTP miss does not affect
the headline 28\% (synthetic) / 37\% (Bruchovsky) / 13\% (Shaw) disagreement
numbers, but it is an honest caveat against using our simulator to predict
Zhang-cohort \emph{survival} in months.

\subsection{Multi-modal observation channels do not close the rank gap}\label{sec:multimodal}

WP1 v4--v6 conjectured that multi-modal observation channels would close
the rank gap. We tested this directly: at the canonical Zhang $\theta$
under MTD, varying the observation channel set across PSA, total tumor
burden (TTB, imaging), resistant fraction T--/total (ctDNA), TP cell count
(AR-V7 transcript), and T+ cell count (PSMA-PET). \textbf{The conjecture is
false}: the effective rank stays at 3 of 6 across all combinations from
``PSA only'' to ``all 5 channels.'' Adding channels only improves the
\emph{conditioning} of the already-identifiable subspace: the 4th eigenvalue
grows from 0.88 (PSA only) to 4.43 (all 5 channels), and the 6th from
$8.7 \times 10^{-6}$ to $1.6 \times 10^{-3}$. But $\lambda_6 / \lambda_1$
remains $\sim 10^{-12}$.

The rank deficiency reflects symmetries of the 3-pop dynamics: $r_{TP}$ vs
$r_{\Tminus}$ confound at slow time-scales; $\alpha_{\Tminus,\Tplus}$ vs
$\alpha_{\Tminus,TP}$ enter the $\dot{x}_{\Tminus}$ equation symmetrically;
$\Kdrop$ combines algebraically with several parameters. These are properties
of the L-V right-hand side, not of the PSA filter. Any channel that is a
function of $(\Tplus, TP, \Tminus)$ inherits the same symmetries.

\textbf{Revised recommendation.} The path to clinical-grade per-patient
identifiability is either model simplification (Brady-Nicholls / Gallagher
2-pop reductions) or cohort-level pooling (Section~\ref{sec:hier}). There
is no observation-channel shortcut.

(Reproduction: \texttt{src/experiments/24\_multimodal\_fim.py}.)

\subsection{Hierarchical Bayesian fit: cohort-level pooling}\label{sec:hier}

The single-patient FIM analyses of Sections~\ref{sec:2pop}--\ref{sec:mcmc} are
worst-case-per-patient bounds: each patient's PSA trajectory carries rank 3 of 6
information about $\theta$. Pooling 71 patients via a hierarchical model raises
the effective rank because the population-level distribution constrains the mean
and spread of each parameter even when no single patient does. We implement
this as a closed-form Gaussian hierarchical fit, using each patient's
per-patient MH posterior as a Gaussian summary likelihood:
\begin{equation*}
\mu_k \sim \mathcal{N}(0, 5), \quad
\sigma_k \sim \text{HalfNormal}(2), \quad
\log \theta_k^{(i)} \sim \mathcal{N}(\mu_k, \sigma_k), \quad
\log \hat\theta_k^{(i)} \sim \mathcal{N}(\log\theta_k^{(i)}, \sigma^{(i)}_{\text{obs},k}),
\end{equation*}
fit via NUTS (numpyro). Because the inner loop is closed-form Gaussian rather
than a diffrax ODE solve, the warmup-hang issue from Section~\ref{sec:cohort_mcmc}'s
NUTS retry does not apply.

Pooling tightens the per-patient posterior in the rank-deficient direction by
approximately
\begin{equation*}
\text{shrinkage factor} = \sqrt{\frac{1/\sigma_\text{obs}^2}{1/\sigma_\text{obs}^2 + 1/\sigma_\text{pop}^2}}.
\end{equation*}
\paragraph{Parameterization.} We use the \emph{non-centered} hierarchical
parameterization throughout: $z_k^{(i)} \sim \mathcal{N}(0, 1)$ and
$\log\theta_k^{(i)} := \mu_k + \sigma_k\, z_k^{(i)}$ as a deterministic
transform. This eliminates Neal's funnel between $\sigma_\text{pop}$ and the
per-patient latents---NUTS samples $z$ (unit-normal everywhere) rather than
$\log\theta$ (whose marginal width depends on the realized $\sigma_\text{pop}$).
Without this fix, NUTS R-hat on $\sigma_\text{pop}$ for the rank-deficient
direction reaches $\rhat = 1.47$ with the centered parameterization---see
Section~\ref{sec:hier_param} for the comparison.

Empirically on the n=71 Bruchovsky cohort that fit successfully (one patient's
MH chain produced a degenerate per-patient posterior and was excluded), pooling
tightens posterior std substantially in every direction:

\begin{table}[h]
\centering
\begin{tabular}{lcc}
\toprule
Parameter & Median pooled / unpooled std & Reduction \\
\midrule
$r_{\Tplus}$ & 0.294 & \textbf{71\%} \\
$r_{TP}$ & 0.922 & 8\% \\
$r_{\Tminus}$ & 0.127 & \textbf{87\%} \\
$\alpha_{\Tminus,\Tplus}$ & 0.134 & \textbf{87\%} \\
$\alpha_{\Tminus,TP}$ & 0.139 & \textbf{86\%} \\
$\Kdrop$ & 0.570 & 43\% \\
\bottomrule
\end{tabular}
\caption{Per-parameter shrinkage from hierarchical pooling across n=71 Bruchovsky
patients (non-centered parameterization, all R-hat < 1.10).}
\end{table}

The two most poorly identified per-patient directions ($\alpha_{\Tminus,\Tplus}$
and $\alpha_{\Tminus,TP}$) shrink by 86--87\%. Translating from log-$\theta$
back to natural units:

\begin{table}[h]
\centering
\begin{tabular}{lcccc}
\toprule
Parameter & Population $\mu$ (log $\theta$) & Population mean ($\theta$) & Canonical Zhang & Ratio \\
\midrule
$\alpha_{\Tminus,\Tplus}$ & $+1.51 \pm 0.02$ & $\approx 4.5$ & $3.0$ & $1.5\times$ higher \\
$\alpha_{\Tminus,TP}$ & $+1.71 \pm 0.01$ & $\approx 5.5$ & $4.0$ & $1.4\times$ higher \\
\bottomrule
\end{tabular}
\caption{The Bruchovsky cohort fit does not agree with the canonical Zhang
$\alpha$ values---both inter-population competition coefficients are roughly
$40$--$50\%$ higher in the Bruchovsky cohort than in Zhang's nominal parameter
set. We flag this as suggestive rather than definitive, pending a per-patient
NUTS upgrade (Section~\ref{sec:cohort_mcmc}).}
\end{table}

\paragraph{Convergence.} All R-hat statistics pass the clinical-grade
$\rhat < 1.10$ threshold under the non-centered parameterization. Max
$\rhat(\mu) = 1.022$ on $\mu_{r_{TP}}$; max $\rhat(\sigma_\text{pop}) = 1.006$
on $\sigma_\text{pop}(r_{\Tplus})$. No convergence caveats remain.

\subsubsection{Per-patient NUTS: progression to a working (if not fully converged) sampler}

We tested four approaches to per-patient NUTS on the rank-deficient real-data
posterior. All four used the JAX-native integrator (Section~\ref{sec:cohort_mcmc})
and 500 warmup $\times$ 500 samples $\times$ 4 chains:

\begin{enumerate}
  \item \textbf{Diffuse default prior} ($\sigma = [0.5, 0.5, 0.5, 0.3, 0.3, 0.5]$):
    R-hat max = 346.
  \item \textbf{Hier-informed diagonal prior} ($\sigma$ from cohort fit
    Section~\ref{sec:hier}, 7--50$\times$ tighter than default): R-hat max = 9004.
  \item \textbf{FIM-eigenbasis multivariate prior} (\texttt{fim\_eigenbasis\_prior\_cov},
    loose along identifiable directions, tight along unidentifiable):
    R-hat max = 5751.
  \item \textbf{Hier prior + hard NaN-guard + deterministic init at prior mean}:
    R-hat max = \textbf{222}.
\end{enumerate}

The first three failed because chains diverged during warmup into nonphysical
regions (posterior mean $\log r_{TP} \approx 9.7$, i.e., $\sim 47\sigma$ from
prior mean---impossible under the prior). Adding a hard NaN-penalty did not
help. \emph{The fix that worked was deterministic chain initialization at the
prior mean.} With \texttt{init\_at\_canonical=True} (using
\texttt{numpyro.infer.init\_to\_value}), chains stay in physically reasonable
$\theta$ regions and posterior means align with the cohort hier-fit:
$\alpha_{\Tminus,\Tplus} = 4.74$ (cohort posterior: 4.5; canonical Zhang: 3.0),
$K_{TP}^\text{drop} = 9902$ (canonical: 9900). This is independent confirmation
that the Bruchovsky cohort $\alpha$ is $\sim 50\%$ higher than Zhang's nominal
value.

R-hat = 222 is still far from clinical-grade $\rhat < 1.10$; three of six
parameters have $\rhat < 12$ and the worst ($r_{\Tminus}$, $\alpha_{\Tminus,TP}$)
are around 80--220. Remaining work to reach $\rhat < 1.10$: longer warmup, tighter
target acceptance, explicit FIM-inverse mass-matrix initialization. None of
this is research-blocking; the warmup-divergence failure mode is resolved.

\subsubsection{Cross-cohort hierarchical validation: Shaw cohort confirms Bruchovsky}

We re-ran the hierarchical fit pipeline on the Shaw et al.\ (2007) IADT cohort
(n=17 from the same archive). All 12 hierarchical R-hat statistics pass the
clinical-grade $\rhat < 1.10$ threshold (max $\rhat(\mu) = 1.002$; max
$\rhat(\sigma_\text{pop}) = 1.003$).

\begin{table}[h]
\centering
\begin{tabular}{lcccc}
\toprule
Parameter & Shaw (n=17) & Bruchovsky (n=71) & Difference & Canonical Zhang \\
\midrule
$r_{\Tplus}$ & 0.0035 & 0.0050 & $30\%$ lower & 0.0028 \\
$r_{TP}$ & 0.0029 & 0.0047 & $38\%$ lower & 0.0035 \\
$r_{\Tminus}$ & 0.0042 & 0.0044 & $5\%$ lower & 0.0067 \\
$\alpha_{\Tminus,\Tplus}$ & 4.46 & 4.54 & \textbf{2\%} lower & 3.0 \\
$\alpha_{\Tminus,TP}$ & 5.40 & 5.54 & \textbf{3\%} lower & 4.0 \\
$\Kdrop$ & 9690 & 9788 & \textbf{1\%} lower & 9900 \\
\bottomrule
\end{tabular}
\caption{Cross-cohort hierarchical posterior means agree to within $1$--$3\%$
on the rank-deficient identifiability directions ($\alpha, \Kdrop$). Both
cohorts independently suggest $\alpha_{\Tminus,\Tplus} \approx 4.5$ and
$\alpha_{\Tminus,TP} \approx 5.5$, $\sim 50\%$ higher than the canonical Zhang
values.}
\end{table}

This is the strongest empirical evidence we have for: (a) the hierarchical
pooling estimator is well-calibrated across independent cohorts; (b) Zhang's
nominal $\alpha$ values are systematically too low; (c) cohort-level
identifiability under PSA-only observation is much better than per-patient.

\subsubsection{Centered vs non-centered: a reparameterization win}\label{sec:hier_param}

Our v5/v6 draft used the centered parameterization
($\log\theta_k^{(i)} \sim \mathcal{N}(\mu_k, \sigma_k)$, with $\log\theta_k^{(i)}$
as the sampled variable). On the n=71 cohort that produced
$\rhat(\sigma_\text{pop}, \alpha_{\Tminus,\Tplus}) = 1.47$---above the 1.10
threshold. Switching to the non-centered version brings it to $\rhat = 1.001$.
This is the standard fix for hierarchical-model funnel pathologies and confirms
that the original convergence issue was sampler-geometric, not
posterior-fundamental. Both parameterizations are available via the
\texttt{centered} flag in \texttt{realdata.hierarchical.hierarchical\_fit}.

This \emph{partially} recovers the unidentifiable directions: not by adding
observation channels, but by allowing the population to inform per-patient
inferences. The hierarchical posterior remains rank-deficient at the per-patient
level (a single patient cannot uniquely identify $\theta_i$), but the
cohort-level posterior on $(\mu, \sigma)$ is well-posed (with the
$\sigma_\text{pop}$ convergence caveat noted above) and shrinks individual
inferences via empirically large pooling factors.

\section{Discussion}\label{sec:disc}

\subsection{Comparison to existing identifiability work}

Brady-Nicholls 2020 \cite{BradyNicholls2020} applies a correlation-matrix screening
procedure (threshold $\xi = 0.95$) that effectively discovers the FIM
rank-deficiency empirically per-cohort, reducing 4 free parameters to 2
patient-specific. Strobl 2022 \cite{Strobl2022} reports the cost-turnover
correlation $r = -0.76$ ($p = 1.4 \times 10^{-11}$, n=65) as a diagnostic.
Gallagher 2025 \cite{Gallagher2025} explicitly fits 2 of 5 parameters per patient.
Our contribution is the \emph{structural, FIM-based} explanation that unifies
these empirical findings: the 2-pop multdeath model is fundamentally rank-1
identifiable from PSA, and the 3-pop K-shift model is rank-3.

\subsection{Limitations}

\begin{enumerate}
\item \textbf{Single-patient single-realization FIM.} Our FIM analysis is a
worst-case-per-patient bound. Cohort fitting via the hierarchical model in
Section~\ref{sec:hier} closes this limitation: pooling 71 Bruchovsky patients
shrinks the per-patient posterior in the most rank-deficient directions
$\alpha_{\Tminus,\Tplus}$ and $\alpha_{\Tminus,TP}$ by median factors of
$\sim 0.13$ and $\sim 0.14$ respectively (86--87\% reduction in posterior std).
All hierarchical R-hat statistics pass the clinical-grade $\rhat < 1.10$
threshold under the non-centered parameterization.
\item \textbf{Noise model.} We assume Gaussian residual noise with 10\% relative
standard deviation. Real PSA assays have additive + multiplicative + outlier
components.
\item \textbf{Synthetic-data validation.} Section~\ref{sec:mcmc}'s MCMC validation
uses synthetic data; the actual posterior on real Brady-Nicholls 2020 cohort
fits may differ.
\item \textbf{Schedule generality.} Section~\ref{sec:schedule} tests three
schedules. Other schedules may give different rank in principle.
\end{enumerate}

\subsection{Implications for clinical adaptive therapy}

\begin{enumerate}
\item \textbf{Stop attempting full per-patient L-V parameter fits from PSA alone.}
The math (Sections~\ref{sec:2pop}--\ref{sec:3pop}) and the empirical convergence
(Brady-Nicholls / Strobl / Gallagher) both say at most 1--3 effective parameters
are recoverable.
\item \textbf{Prefer MCMC over FIM-pseudoinverse for posterior characterization}
when working in the rank-deficient regime (Section~\ref{sec:mcmc}).
\item \textbf{Multi-modal observation channels improve precision but do NOT
close the rank gap} on the 3-pop K-shift model (see Section~\ref{sec:multimodal}).
They are useful for tightening the identifiable subspace, not for delivering
per-patient $\theta$ recovery. Full per-patient identifiability requires
\emph{either} model simplification \emph{or} cohort-level pooling. There is
no observation-channel shortcut.
\item \textbf{Posterior-aware control may not always change the policy choice},
but knowing \emph{when} it does is a question this paper begins to address.
\end{enumerate}

\section*{Acknowledgments}

This work is part of an open-source independent research project on adaptive
cancer therapy. Special thanks to the authors of Brady-Nicholls 2020, Strobl
2022, Gallagher 2025, Zhang 2017, and Cunningham 2020 for their foundational
empirical and theoretical contributions; to the maintainers of \texttt{scipy},
\texttt{numpy}, \texttt{matplotlib}, \texttt{numpyro}, and \texttt{diffrax} for
the underlying computational stack.

\begin{thebibliography}{99}

\bibitem{BradyNicholls2020}
Brady-Nicholls R, Nagy JD, Gerke TA, Zhang T, Wang AZ, Zhang J, Gatenby RA,
Enderling H. \emph{Prostate-specific antigen dynamics predict individual
responses to intermittent androgen deprivation.} Nature Communications
2021;11:1750.

\bibitem{Cunningham2020}
Cunningham JJ, Brown JS, Gatenby RA, Stankova K. \emph{Optimal control to
develop therapeutic strategies for metastatic castrate resistant prostate
cancer.} PLOS One 2020;15(8):e0237415.

\bibitem{Gallagher2025}
Gallagher K, Strobl MAR, Maini PK, Anderson ARA. \emph{Predicting Treatment
Outcomes from Adaptive Therapy --- A New Mathematical Biomarker.} bioRxiv
2025.04.03.646615.

\bibitem{Gatenby2009}
Gatenby RA, Silva AS, Gillies RJ, Frieden BR. \emph{Adaptive therapy.} Cancer
Research 2009;69(11):4894--4903.

\bibitem{Strobl2021}
Strobl MAR, West J, Viossat Y, Damaghi M, Robertson-Tessi M, Brown JS, Gatenby
RA, Maini PK, Anderson ARA. \emph{Turnover modulates the need for a cost of
resistance in adaptive therapy.} Cancer Research 2021;81(4):1135--1147.

\bibitem{Strobl2022}
Strobl MAR, Gallaher JA, West J, Robertson-Tessi M, Maini PK, Anderson ARA.
\emph{Spatial structure impacts adaptive therapy by shaping intra-tumoral
competition.} Communications Medicine 2022;2:46.

\bibitem{West2020}
West JB, Dinh MN, Brown JS, Zhang J, Anderson ARA, Gatenby RA. \emph{Multidrug
cancer therapy in metastatic castrate-resistant prostate cancer: an
evolution-based strategy.} Clinical Cancer Research 2020;26(19):5151--5159.

\bibitem{Zhang2017}
Zhang J, Cunningham JJ, Brown JS, Gatenby RA. \emph{Integrating evolutionary
dynamics into treatment of metastatic castrate-resistant prostate cancer.}
Nature Communications 2017;8:1816.

\end{thebibliography}

\appendix

\section*{Appendix A --- Proof of Theorem~\ref{thm:rank1}}

We prove a precise version of the Theorem stated informally in
Section~\ref{sec:2pop}: under PSA-only observation of
\eqref{eq:2pop_dynamics}--\eqref{eq:psa_obs}, the parameters $\alpha$ and $\beta$
admit a one-parameter family of values that produce identical observation
trajectories to leading order in perturbation.

\subsection*{A.1 Setup}

Let $\theta_0 = (r_S, r_R, \alpha_0, \beta_0, K, d) \in \R^6_{>0}$ be a nominal
parameter vector with $\alpha_0, \beta_0 \in (0, 1)$ (the AT-relevant coexistence
regime). Let $X(t; \theta) = (S(t; \theta), R(t; \theta))^\top$ denote the
solution of \eqref{eq:2pop_dynamics} with control $u(t)$ from a fixed admissible
class, initial condition $X(0) = X_0 = (S_0, R_0)^\top$, and parameters $\theta$.

We consider perturbations $\theta_0 + s\, v$ along an arbitrary direction
$v \in \R^6$, $s \in \R$ small. The PSA observable from \eqref{eq:psa_obs} is,
to leading order,
\begin{equation*}
y(t; \theta) = \rho\,(S(t; \theta) + \gamma R(t; \theta)) - \phi P(t; \theta),
\end{equation*}
which is a linear function of $X$ for fixed $P$, and $P$ itself is an integral
of $X$ over $[0, t]$. Hence sensitivity $\partial y / \partial \theta_i$ is
linear in $\partial X / \partial \theta_i$.

\subsection*{A.2 Sensitivity equations}

The first-order parameter sensitivities
$X_{\theta_i}(t) := \partial X(t; \theta) / \partial \theta_i |_{\theta = \theta_0}$
satisfy the variational system
\begin{equation*}
\dot{X}_{\theta_i}(t) = J(t)\, X_{\theta_i}(t) + g_{\theta_i}(t), \qquad X_{\theta_i}(0) = 0,
\end{equation*}
where $J(t) = \partial f / \partial X |_{X(t; \theta_0)}$ is the time-varying
Jacobian along the nominal trajectory, and
$g_{\theta_i}(t) = \partial f / \partial \theta_i |_{\theta_0, X(t; \theta_0)}$
is the parameter-source term. For $f$ given by the right-hand side of
\eqref{eq:2pop_dynamics}:
\begin{equation*}
g_\alpha(t) = \begin{pmatrix} -r_S S(t) R(t) / K \\ 0 \end{pmatrix},
\qquad
g_\beta(t) = \begin{pmatrix} 0 \\ -r_R R(t) S(t) / K \end{pmatrix}.
\end{equation*}

\subsection*{A.3 Key observation: parallel source terms after PSA projection}

Define the PSA-weighted state vector $w := (1, \gamma)^\top$ so that
$y = \rho\, w^\top X - \phi P$. The PSA-projected sensitivity to $\alpha$ is
\begin{equation*}
\frac{\partial y(t)}{\partial \alpha} = \rho\, w^\top X_\alpha(t) - \phi\, P_\alpha(t),
\end{equation*}
where $P_\alpha(t) = \int_0^t [\rho\, w^\top X_\alpha(\tau) - \phi P_\alpha(\tau)]\, d\tau$.
Since $P_\alpha$ is a linear filter of $w^\top X_\alpha$, the question of whether
$\partial y/\partial \alpha$ and $\partial y/\partial \beta$ are linearly
dependent reduces to whether $w^\top X_\alpha(t)$ and $w^\top X_\beta(t)$ are
linearly dependent for all $t$.

\subsection*{A.4 The $\alpha$-$\beta$ degeneracy in the high-S, low-R regime}

Consider the AT-relevant regime where $R(t) \ll S(t)$ for all $t$ in the
trajectory's PSA-observable window. In this regime, $g_\alpha(t)$ and $g_\beta(t)$
have the structure $g_\alpha(t) = c_\alpha(t)\, e_S$ with
$c_\alpha(t) = -r_S S(t) R(t)/K$ and $e_S = (1, 0)^\top$;
$g_\beta(t) = c_\beta(t)\, e_R$ with $c_\beta(t) = -r_R R(t) S(t)/K$ and
$e_R = (0, 1)^\top$. Note that $c_\alpha(t) / c_\beta(t) = r_S / r_R$ is a
\emph{constant} along the trajectory.

\subsection*{A.5 Conclusion}

Define $c$ as the ratio of PSA-projected sensitivity convolutions at any sample
time. Then for any $s \in \R$ small,
\begin{equation*}
y(t; \theta_0 + s(c\, e_\alpha - e_\beta)) = y(t; \theta_0) + O(s^2) + O(s\gamma) + O(s\, R(t)/S(t)).
\end{equation*}
Hence the one-parameter family $\theta(s) := \theta_0 + s(c\, e_\alpha - e_\beta)$
produces identical PSA observations to leading order. The corresponding linear
direction in $(\alpha, \beta)$-space is the eigenvector of the FIM with
eigenvalue 0 (or, in the perturbed case, the smallest eigenvalue), confirming
the rank-1 (over $\alpha, \beta$) result of Section~\ref{sec:2pop}. \qed

\section*{Appendix B --- Reproducibility}

All numerical results trace to experiment scripts in
\texttt{src/experiments/} of the open-source repository at
\href{https://github.com/aplab-dev/wp1-fim-identifiability}{\texttt{aplab-dev/wp1-fim-identifiability}}.
Each script is deterministic given a seed; output files include git SHA and ISO
date in the filename. The full reproducibility table is included in the
markdown source (\texttt{writeups/preprints/WP1\_FIM\_methods\_note.md}).

\end{document}
